\section{The cohomological bound}\label{sec:cohomology}

Item~17 of the note's timeline claims ``a significant weakening of Vandiver'' as a by-product.
This section says what the formal statement is, where the slope argument uses it, and how it is
proved.

\subsection{The statement}
Let $E$ be a number field, $\ell$ a prime, and $e\ge1$. Let $S$ be the set of places of $E$
above $\ell$ and $\infty$, and $G_S$ the Galois group of the maximal extension of $E$ unramified
outside $S$. Define \lvS
\[
\Sha(E,\ell,e)=\ker\Bigl(H^1\bigl(G_S,\Q_\ell/\Z_\ell(-e)\bigr)\longrightarrow
\prod_{v\mid\ell}H^1\bigl(E_v,\Q_\ell/\Z_\ell(-e)\bigr)\Bigr),
\]
the classes that are locally trivial above $\ell$, and let $\len_\ell(e)$ be its length. By
duality this is governed by $H^2(G_S,\Z_\ell(e+1))$, and hence, through Soul\'e's Chern class
maps \cite{soule1979}, by the even $K$-groups $K_{2e}(\cO_E)\otimes\Z_\ell$. For $E=\Q$ the
vanishing of $K_{4i}(\Z)\otimes\Z_\ell$ for all $i$ is equivalent to Vandiver's conjecture at
$\ell$ \cite{kurihara1992}; over $\Q(\mu_\ell)$ the conjecture concerns the eigenspaces of one
parity \cite{washington1997}.

\begin{statement}[\lean{VVle}, \lean{VWin}]\label{st:vwin}
For the host field $E$ at scale $x$, with $t=\lfloor x^{401}\rfloor$ twists, every
$\len_\ell(e)$ with $1\le e<t$ is finite and
\[
\sum_{\ell}\log\ell\sum_{1\le e<t}\len_\ell(e)\ \le\ \sqrt{b},\qquad b=\exp(x^{\eps}),
\]
for all $x$ in a window $[W,W^2]$; there is $J\ge1$ such that for every large $X$ such a window
lies in $[X^{1/J},X]$.
\end{statement}
Nothing is averaged or discarded: every prime $\ell$ and every twist is bounded. The statement
is far from Vandiver's conjecture, which would make part of these groups vanish for
$\Q(\mu_\ell)$; it says that the total size of these groups, for a field of degree about $x^2$
and conductor exponential in a power of $x$, is at most exponential in a small power of $x$.

\subsection{Where it is used}
The bound has a single use \lvP. The record of \cref{sec:cast} carries, for each lift
level $j$, a quantity $\eta_j$: the normalised logarithm of the index of a sublattice
$\mathrm{cut}_j$ in a lattice $T_j$, summed over the places above the level primes. The clause
\lean{LinkN_at} bounds it by a size term plus the total of \cref{st:vwin},
\[
\eta_j\ \le\ d_j\Bigl(C\bigl(fmb\,L^2+(\log x)^2N_c\bigr)+\sum_\ell\log\ell\sum_{1\le e<t}\len_\ell(e)\Bigr),
\]
and the consumer needs only $\eta_j\le D_0\,C\,x^{\eps}(B_0+N_c)$ with $B_0=fmb\,x^\eps$. The
lemma that passes from one to the other (\lean{divinN_of_hypV}) weakens $\sqrt b$ to $b$ at
once. \emph{So the slope argument needs from \cref{st:vwin} only that the total is at most
$\exp(x^{\eps})$}; the square root is slack. The Sha lengths in the record are the true ones:
they are tied to the Galois-cohomological definition by a field of the base data
(\lean{vvRecord}). In the source the weakened bound is used in one place besides its producer
(\lean{inEta_num}).

What it controls is a congruence module \lvS. At each place $w$ above a prime $\ell$ the record
has a local Hecke algebra $\mathbf T$, in spectral form: a subalgebra of the functions from a
finite set $X$ of eigensystems to a discrete valuation ring $\cO$. A subset $Z\subset X$
consists of the eigensystems that have an automorphic point whose parameter contains a block
of length $t$ (\lean{ZN}); these are the lifts. The \emph{congruence ideal} is the kernel of
$\mathbf T\to\prod_{y\notin Z}\cO$ (\lean{ccIdeal}), and the local part of $\eta_j$ is the
length of the corresponding congruence module, $V_Z\cap(L+V_{Z^c})$ modulo $L\cap V_Z$ for
the lattice $L$ of forms (\lean{etaLoc}).

\begin{statement}[\lean{CongBound_at}]\label{st:cong}
Let $\ell$ be large and not exceptional, $\pi\in Z$, and suppose that the Galois
representation of $\pi$ is, over the fraction field, the sum of a representation $\rho$ of
rank $r$ and the $t$ characters $\chi_1,\dots,\chi_t$. Then the valuation at $\pi$ of the
congruence ideal is at most
\[
\len\,\mathrm{Sel}(\rho)\;+\;\sum_{0<e<t}\len_\ell(e)\;+\;v_\ell(h_E)\;+\;t\sum_{P\in Q}v(\zeta_P-1),
\]
where $\mathrm{Sel}(\rho)$ is a Selmer group attached to $\rho$, $h_E$ is the class number of
the host field and $Q$ is a finite set of auxiliary primes.
\end{statement}
This is a bound of the type of Ribet's lemma \cite{ribet1976}: a congruence between a lift and
a form that is not a lift produces a cocycle, and cocycles are counted by Selmer groups, by the
groups $\Sha(E,\ell,e)$ and by the class group. The twists $1\le e\le t-1$ are those of the
block: in the source the characters are $\chi_c=\chi_1\,\chi_{\mathrm{cyc}}^{-(c-1)}$, with
$\chi_{\mathrm{cyc}}$ the cyclotomic character.

This step uses Galois representations attached to every eigensystem, lifted or not. They are
supplied by a theorem of the archive (\cref{sec:galois}).

\subsection{How it is proved}
One theorem assembles \cref{st:vwin} \lvS{} (\lean{of_four_alt}). Its hypotheses are a record
of statements from the literature, the crude bound \lean{LemmaK} below, and four counting
statements (\lean{CountI}, \lean{CountII}, \lean{CountIII}, \lean{CountV}), each of which
bounds the contribution of one class of pairs $(\ell,e)$; its proof calls, besides these, the
theorem of Ferrero and Washington and the bound for the odd part.

\emph{The odd part.} For a character $\psi$ of $\Gal(E/\Q)$ whose value at complex conjugation
is $(-1)^{e+1}$, the length of the $\psi$-eigenspace of the group with twist $e$ is at most the
valuation of an explicit number attached to $\psi$ and $e$ (\lean{PsiBound}). This is proved
by a theorem \lvS{} (\lean{ML.iv}) whose hypotheses are Poitou--Tate duality and an
orthogonality statement for unramified classes. The archive's proof contains a step of
Kolyvagin's derivative construction (\lean{kolStep}); compare
\cite{kolyvagin1990,rubin2000}. The eigenspaces of the other parity are the side on which
Vandiver's conjecture bears.

\emph{Cyclotomic towers.} The theorem of Ferrero and Washington
\cite{ferrerowashington1979} that $\mu=0$ is proved in the archive, by Sinnott's method
\cite{sinnott1984} \lvS{} (\lean{sinnott}); compare \cite{iwasawa1973}. The twists it does not
reach are bounded crudely \lvS{} (\lean{LemmaK}): for a totally complex field
$L\supset\mu_\ell$ and every $e\ge1$, the group $H^2(G_S,\Z_\ell(e+1))$ is finite and the
logarithm of its order is at most an explicit function of $e$, $[L:\Q]$ and $\abs{d_L}$;
compare \cite{soule1979,borel1974}.

\emph{One integer for all $\ell$.} For $\U(3)$ over an imaginary quadratic field there is a
statement of the following shape \lvS{}
(\lean{NumberU3U_s}): for $K$ imaginary quadratic there are constants $c,A$ such
that for every twist $e$, level $N$, prime $p\ge7$ not dividing $N\,d_K$ and Galois-stable set
of roots of unity $z$, there is a non-zero integer $NN$ with
\[
\log\abs{NN}\le c\cdot\#\{z\}\cdot\bigl(N(e+1)\bigr)^A\log p
\]
such that, for every $\ell\ne2,p$, every finite set of pairwise non-isomorphic admissible
Galois representations $\rho$ of that weight and level satisfies
\[
\abs{NN}_\ell\ \le\ \prod_{\rho}\bigl|\tr\rho(\sigma^{-1})-\mathrm{eis}\bigr|_\ell,
\qquad \mathrm{eis}=(\mathrm Nw)^{-e}z+1+(\mathrm Nw)^{e}z^{-1},
\]
at a Frobenius element $\sigma$ at a place $w$ above $p$, with every factor non-zero. Here $\mathrm{eis}$
is the trace at $\sigma$ of a reducible representation $\Psi\oplus1\oplus\Psi^{-1}$, so $NN$ is
a congruence number between cusp forms and an Eisenstein class. Summing $\log\ell$ times the
valuations over all $\ell$ therefore costs only $\log\abs{NN}$.

\emph{From the integer to Sha.} A companion statement \lvS{} (\lean{OneCharU_s}) bounds the
length in question by $2\sum_{\rho}v_\ell\bigl(\tr\rho(\sigma)-t_\Psi\bigr)$ over such a set of
representations, where $t_\Psi=\Psi(\sigma)+1+\Psi(\sigma)^{-1}$, provided no admissible $\rho$ has
trace exactly $t_\Psi$. Behind it is a lemma of commutative algebra \lvS{} (\lean{LemmaM}) about a
graded complete local algebra with an augmentation: under five conditions, the length of the
weight $-1$ part of its cotangent space is at most $2\sum_\xi v(\xi(a))$ over a finite covering
set of points. This is the pattern in which an Eisenstein congruence bounds a Selmer group, as
in the work of Ribet \cite{ribet1976}, Mazur and Wiles \cite{mazurwiles1984} and Skinner and
Urban \cite{skinnerurban2014}. Two further theorems that the archive proves for this purpose
are on the path of the final theorem \lvS: one on Hida families for unitary groups
(\lean{hida_unitary_of_printed_hida}; compare \cite{hida1986}), and Kato's Euler system
\cite{kato2004} (\lean{zetaUr}); compare \cite{blochkato1990}.

No statement of the main conjecture of Iwasawa theory appears among the results the archive
takes from the literature (\cref{sec:literature}).

\subsection{Galois representations from theta lifts}\label{sec:galois}
The representations needed above are supplied by one theorem of the archive \lvS{}
(\lean{galoisRepresentations_of_thetaLift}): for a CM field $E$, a prime $\ell$ and a
regular algebraic, polarised, cuspidal automorphic representation $\pi$ of $\GL_a$ over $E$
there is a continuous semisimple $\ell$-adic representation of $\Gal(\overline E/E)$ with the
Hodge--Tate weights prescribed by $\pi$ and with stated properties at the places above $\ell$.
In the literature this is the work of Clozel, Harris--Taylor, Shin and Chenevier--Harris
\cite{clozel1991,harristaylor2001,shin2011,chenevierharris2013}. Two features of the formal
version should be recorded.

First, its compatibility with $\pi$ away from $\ell$ is partial \lvS{} (\lean{AwayB}): the
dimensions agree; the parameters are equal when $\pi$'s is a sum of characters with trivial
monodromy; and for tame, square-free parameters at places with completion $\Q_g$, $g>16a^2$,
the characteristic polynomials on inertia agree. This is less than the theorem of Caraiani
\cite{caraiani2012}.

Second, the theorem has a single hypothesis, which the archive proves \lvS{}
(\lean{TLiftK}, proved as \lean{tLiftK}): theta lifts exist between \emph{definite} unitary
groups. Precisely, for a Hermitian space of rank $n$ and a skew-Hermitian space of rank
$d\ge d_0$ over $E$, both definite, and an automorphic representation of $\U(n)$ whose weights
satisfy an inequality against $d$, there is a lift to $\U(d)$; compare the work of J.-S.~Li
\cite{li1992}. Neither an eigenvariety nor a stable trace formula appears in the statement.
